\section{Conclusiones}
\label{iii:sec:conclusiones}
\begingroup
\fontsize{10.5}{12.8}\selectfont
\setlength{\parskip}{2pt plus .5pt minus .5pt}

El resultado central es una determinación constitutiva reversible en
el dominio contractivo: los dos canales orientados de un transporte
producen cuatro coordenadas relacionadas y esas coordenadas permiten
recuperar el par angular. La estructura fija primero la ecuación de
respuesta; su evaluación determina los valores sin perder la orientación
necesaria para reconstruir sus antecedentes espectrales. Esta doble
dirección, constructiva e inversa, da contenido preciso a la tesis
estructura, ecuación y valor.

\subsection*{La respuesta conserva el par angular}

Las incidencias $90$ y $120$ determinan una única solución de la
ecuación de completación para el transporte fijado. Su función escalar
es una biyección estrictamente decreciente de $(0,1)$ sobre $(3/4,1)$.
Por ello, la ecuación cúbica~\eqref{iii:eq:cubic} posee exactamente
una raíz admisible para cada respuesta de canal. La inversión no es
una coincidencia numérica: reconstruye las dos contracciones y, por
sus logaritmos, el par angular.

Las cuatro coordenadas del vacío satisfacen
\[
\widehat\mu\widehat\varepsilon=\widehat c^{-2},\qquad
\widehat\mu/\widehat\varepsilon=\widehat Z^2.
\]
Dentro de las reglas de producto, cociente y positividad fijadas, estas
identidades determinan las evaluaciones cuadráticas. Su dependencia
común se manifiesta también bajo intercambio de hojas: permitividad y
permeabilidad se permutan, impedancia se invierte y velocidad se
conserva. El producto simétrico y el cociente orientado transmiten así
informaciones complementarias. La recuperación mantiene las etiquetas
de hoja y la normalización interna; el estado de origen conserva
además ruta, frontera y memoria.

La factorización $\mathcal V=a\exp(\eta\mathcal R)$ identifica
esas dos informaciones en el propio operador. El factor positivo
$a=\sqrt{\det\mathcal V}$ determina la velocidad normalizada;
la componente unimodular determina la impedancia. Conservar ambos
permite reconstruir la respuesta. Esta separación explica la relación
entre las hojas recíprocas de producto unitario y las cuatro lecturas
del vacío: la normalización separa el modo común, que permanece
necesario para recuperar el operador completo.

La inversa permite enlazar respuestas que evalúan funciones distintas
de los mismos canales. La razón de sus exponentes reconstruye la
anisotropía de la elipse de acción y la impedancia LC relativa,
según la proposición~\ref{iii:prop:forma-LC}. La sección de acción
y las referencias de frecuencia e impedancia completan sus magnitudes dimensionales.
La identidad $c_{\rm int}=c_{\rm vac}$ establece, a su vez, que el
lector cinemático y el constitutivo realizan la misma sección de
velocidad en una coordenatización común. La traza normalizada conserva esa
lectura bajo amplificación nonádica.

\subsection*{Composición y reconstrucción entre sectores}

La ley transportada $r\odot t=f(\psi(r)\psi(t))$ expresa la composición
del transporte en coordenadas constitutivas. Su neutro es $3/4$ en la
extensión de frontera, y $-\log\psi(r)/30$ convierte esa composición en
suma. Los proyectores comunes preservan la realización operatoria y las
etiquetas de hoja. La velocidad de un único estado conserva, por su parte,
el producto ordinario de sus dos canales. Ambas operaciones quedan
distinguidas por sus argumentos y por la prueba exacta de su relación.

La reconstrucción Higgs--impedancia establece otra conexión explícita.
Los acoplamientos de calibre recuperan la contracción común y el
acoplamiento electromagnético; el autoacoplamiento de Higgs recupera la
coordenada orientada. En el dominio de la realización de árbol, esos
datos determinan los dos canales y sus respuestas del vacío. Las
desigualdades de dominio y las compatibilidades genealógicas precisan
qué ternas pertenecen a la realización considerada. Conservar los
proyectores marcados permite restituir también el operador; la historia
enriquecida mantiene sus datos de procedencia adicionales.

\subsection*{El ciclo relaciona respuesta, momento y orientación}

La realización dodecafásica explica la procedencia de la función
racional por un cociente de determinantes en dimensiones $3$ y $4$.
La diferencia de rangos produce $-1/12$ como traza normalizada en ese
dominio finito. El plano de cuarto de giro conserva otra información:
su operador satisface $J^2=-I$ y su coeficiente de resolvente retiene
el signo de orientación. Las dos integraciones logarítmicas de ese
coeficiente, con sus condiciones iniciales, determinan el momento
de Catalán a profundidad ilimitada.

La tabla de trazas $b_n$ reconstruye $-\log f(s)$ y produce momentos
convergentes para $\operatorname{Re}z>0$, con una cota explícita del
resto. Su valor en profundidad uno es $\log(4/3)$, compatible con
el extremo constitutivo. La representación de Mellin y la distinción
entre producto diagonal y tensorial conservan las operaciones que
preceden a la evaluación escalar de cada momento.

La restitución integral del
teorema~\ref{iii:thm:restitucion-funcional-catalan} cierra la
correspondencia funcional en ambas direcciones:
\[
 \mathcal C_2(s)=\int_0^s\log\frac{s}{t}\,
       \frac{1+t}{1+t+t^2}f(t)\,dt,\qquad
 f(s)=\frac{1+s+s^2}{s(1+s)}\mathcal D^2\mathcal C_2(s).
\]
La condición $\mathcal C_2(s)\to0$ en el origen elimina toda
libertad aditiva y logarítmica. El valor extremo $G_{\rm Cat}$,
el núcleo orientado y la respuesta constitutiva quedan así situados
en una misma colección, con operaciones de restitución explícitas.
Los argumentos $s_\pm$ conservan su determinación angular; la traza
normalizada mantiene la velocidad del mismo estado al enlazarla
con la acción y la longitud. Esta correspondencia especifica el
contenido estructural que se conserva al pasar de la colección a sus valores.

La identidad Catalán--Pell precisa una evaluación algebraica de esa
colección:
\[
\mathcal C_2(2-\sqrt3)
=\frac23G_{\rm Cat}-\frac{\pi}{12}\log(2+\sqrt3).
\]
El coeficiente $2/3$ procede de las clases impares del calendario y
del peso de profundidad; las dos primeras derivadas logarítmicas en esa sección toman los valores
$\pi/12$ y $1/4$. La reciprocidad de la unidad de Pell enlaza
expansión y contracción, mientras los dos argumentos constitutivos
siguen determinados por el par angular. Esta relación mantiene
visibles la colección, sus lectores y las secciones en que se evalúan.

\Needspace{9\baselineskip}
El funcional de Barbero--Immirzi se expresa después como diferencia
orientada de una misma función de los canales reconstruidos:
\[
\gamma=\mathcal H(\psi(r_-))-\mathcal H(\psi(r_+)).
\]
La incidencia y el retorno dodecafásico fijan sus términos; la
inversión constitutiva permite transportarlos a la respuesta del
vacío. Con la marca de orientación fija, el intercambio de canales
cambia el signo de $\gamma$; la conjugación conjunta conserva el
funcional. Esa distinción prueba que su contenido excede el espectro
sin etiquetas y se mantiene bajo las amplificaciones normalizadas.

% BEGIN SINTESIS_ADITIVA_20260922
\par
La sección~\ref{delta22:iii:restitucion} reúne el potencial espectral de la respuesta constitutiva, la restitución global de los canales mediante velocidad normalizada y Barbero--Immirzi, y sus cotas de estabilidad. Sobre la imagen generada, con las etiquetas y bases conservadas, la composición recupera el par angular, la sección de acción y la lectura periódica del registro dodecafásico.
% END SINTESIS_ADITIVA_20260922

La resolución armónica del patrón dodecafásico conserva exactamente
los modos $3,9,4,8$, con coeficientes no normalizados $72,72,-6,-6$.
La razón de módulos $12:1$ expresa las multiplicidades orientadas.
El defecto de incidencia $-1/12$ y el coeficiente singular $1/24$
siguen ligados a sus propias operaciones. La comparación explícita
con el lector de pesos iguales, cuyo coeficiente es $1/1080$,
precisa cómo interviene la cadena orientada en la normalización
del retorno sin sustituir los funcionales por sus polos.

\subsection*{Del balance aritmético a las relaciones eléctricas}

La discrepancia de vacancias aporta una polarización centrada cuya
magnitud permanece menor que una unidad y cuyo incremento satisface
un balance temporal exacto. Los cambios entre las convenciones de
suelo y techo quedan registrados por un término de frontera. Una
sección de carga y una duración orientada realizan ese balance como
carga centrada y corriente, conservando su ley telescópica.

El prefactor electrónico conserva asimismo su estructura antecedente.
La enumeración de las fibras en $D_9^3$ fija el modo singular trítico
de valor $1/2$; el espectro completo de su emparejamiento prueba
\[
\|[P_{\Sigma,3},P_{\Pi,3}]\|=\sqrt3/4.
\]
Su defecto cuadrático $3/4$ coincide con el límite del cociente de
sectores $90/120$. Esta igualdad relaciona resultados de operadores
especificados, cada uno con su espacio y su acción: el prefactor es
una norma de incompatibilidad aritmética, mientras la respuesta
interior depende de sus canales contractivos.

La sección positiva de carga determinada por
$Z_{\rm vac}e_a^2=2\alpha h_a$ reúne el acoplamiento, la acción y
la impedancia. Las relaciones de von Klitzing, conductancia, flujo
y Josephson satisfacen exactamente
\[
R_KG_0=2,\qquad G_0Z_{\rm vac}=4\alpha,\qquad
K_{J,a}\Phi_{0,a}=1,\qquad K_{J,a}^2R_K=4/h_a.
\]
El retorno de la acción conserva resistencia y conductancia; carga y
flujo se transforman con $R_{\rm act}^{1/2}$ y Josephson con
$R_{\rm act}^{-1/2}$. Estas dependencias explican conjuntamente
qué magnitudes retienen la acción y cuáles conservan únicamente su
cociente constitutivo.

La realización dimensional preserva todas estas identidades al
cambiar las bases de acción, carga y velocidad. Los valores
adimensionales obtenidos pertenecen a la coordenatización interna de la
construcción; su reconocimiento metrológico exige la evaluación
independiente de las secciones dimensionales y del régimen físico
correspondiente. El alcance matemático establecido es la respuesta
positiva e invertible, sus composiciones y su covariancia en los
dominios declarados.

El conjunto muestra cómo una misma genealogía APP--TRIT--TPK conserva
las relaciones necesarias para pasar de caminos y operaciones
aritméticas a canales, momentos, respuesta constitutiva y relaciones
eléctricas. Los valores adquieren así una procedencia explicada y
una información recuperable: expresan las operaciones que los
producen y permiten reconocer, dentro de su representación, la
estructura que permanece en ellos.
\par\endgroup
